\documentclass[11pt,a4paper]{article}
\usepackage[T1]{fontenc}
\usepackage{lmodern,amsmath,amssymb,amsthm,mathtools,booktabs,microtype}
\usepackage[margin=26mm]{geometry}
\usepackage[colorlinks=true,linkcolor=blue,citecolor=blue,urlcolor=blue]{hyperref}
\setlength{\emergencystretch}{2em}
\newtheorem{theorem}{Theorem}[section]
\newtheorem{lemma}[theorem]{Lemma}
\newtheorem{proposition}[theorem]{Proposition}
\newtheorem{corollary}[theorem]{Corollary}
\theoremstyle{remark}\newtheorem{remark}[theorem]{Remark}
\newcommand{\Q}{\mathbb Q}\newcommand{\Z}{\mathbb Z}
\DeclareMathOperator{\Tr}{Tr}\DeclareMathOperator{\rad}{rad}
\title{A quartic in 78 variables from a mixed-power height estimate}
\author{}\date{15 November 2026}
\begin{document}\maketitle
\begin{abstract}
Under ordinary rational $abc$ and the precisely retained arithmetic
inputs of the 79-variable construction, we remove one rational coordinate.
The integer compared with an elliptic index is $X^2+3$, while the sufficiently
large fixed shift is moved into the comparison scale $B+X^8+T^2$.
A parity-sensitive consequence of $abc$ controls both $h(X)$ and $h(T)$
by the height of this scale. The eight-contact estimate gives the strict
absorption coefficient $602/679$. The resulting polynomial is a sum of
70 primitive quadratic squares in 78 total variables, has degree exactly
four and integer content one. Its defining shift remains ineffective.
No matching of the two pole square classes is required.
\end{abstract}

\section{Statement, counting convention, and hypotheses}
Every variable count includes the input. For $x=u/v\in\Q$ in lowest terms,
$v>0$, put $h(x)=\log\max(|u|,v)$; in particular $h(0)=0$.
Put $K=\Q(\rho)$, $\rho^2=2$, $D=11-6\rho$, and
\[
 E:y^2=x^3-Dx,\qquad R_0=(1+\rho,2\rho),\qquad P=2R_0.
\]
Retain exactly the following inputs and applications of the preceding
79-variable source~\cite{Parent}: $2E(K)=\Z P$, with $P$ non-torsion;
the fixed-data simultaneous-prime theorem with its prescribed local
orders, finite avoidance, square denominator, and real approximation;
and the prepared congruent-number pole realization supplied there by
Jaswal's simultaneous-twist theorem~\cite{Jaswal,Pole}.
Ordinary rational $abc$ is the additional conjecture. No new prime or
pole realization hypothesis is introduced. The source's complete
reference-chart and reciprocal-point identities, local extraction
lemma, and all-point trace-height theorem are used with their stated
domains. Their relevant conclusions are restated below. Neither finite
coefficient checks nor this paper constitute an independent audit of the
external arithmetic premises.

\begin{theorem}\label{thm:main}
Under these inputs, there is a computable family
\[
 Q_B\in\Z[X,Y_1,\ldots,Y_{77}],\qquad B\ge81\text{ a fixed odd integer},
\]
such that every member has total degree exactly four and integer content
one. It is a sum of 70 primitive integral quadratic squares. There is a
fixed $B_*$ such that every odd $B\ge B_*$ has rational input projection
exactly $\Z$. The shift $B$ is a coefficient, not a quantified variable.
No numerical sufficient value of $B_*$ is given.
\end{theorem}

The source formulas in Section~\ref{sec:source}, together with the
unchanged, explicitly pinned reference coefficients of~\cite{Parent},
are a finite coefficient specification. The directory
\nolinkurl{computation/quartic/abc78_v34} contains the full quadratic and
quartic tables, the generator, and a separate named-row reconstruction.
A $B=81$ substitution is an algebraic test member, not a certified defining
member. The general counts 25 and 13, and the inherited-premise quartic
count 200, are not reduced by this argument.

\section{Eight contacts under ordinary rational $abc$}
We give the height estimate used here, rather than treating a stronger
contact coefficient as an unspecified assumption.

\begin{lemma}\label{lem:marked}
Let $\theta_1,\ldots,\theta_m$ be distinct rational numbers. Outside a
fixed finite set $S$ of primes containing their denominators and pairwise
differences, let $M(t)$ be the product of primes at which some
$t-\theta_i$ has positive odd valuation. Ordinary rational $abc$ implies,
for every $\delta>0$ and $m>4+2\delta$,
\[
 \log M(t)\ge(m-4-2\delta)h(t)-C_\delta
 \qquad(t\in\Q\setminus\{\theta_1,\ldots,\theta_m\}),
\]
where $C_\delta$ depends on the fixed contacts and $S$, not on $t$.
\end{lemma}
\begin{proof}
We recall the marked Belyi argument, the standard consequence of $abc$
used in~\cite{Elkies}. There is a rational Belyi map
$\varphi:\mathbf P^1\to\mathbf P^1$, of some fixed degree $d$, whose
inverse image of $\{0,1,\infty\}$ contains the marked rational contacts.
The map can be defined over $\Q$: the marked set is rational and the
usual Belyi construction uses a Galois-stable finite set of critical
values. Riemann--Hurwitz gives exactly $d+2$ distinct geometric points
in those three fibers. The remaining divisor therefore has degree
$d+2-m$.

For a primitive pair $(u,v)$, the prime support of the coprime integer
triple representing $\varphi(u/v)$ and $1-\varphi(u/v)$ is, outside
fixed bad primes, contained in the zeros of the marked linear factors
and of that remaining divisor. A homogeneous equation for the remaining
divisor has fixed coefficients and degree $d+2-m$, so its nonzero value
has logarithmic absolute value at most $(d+2-m)h(u/v)+O(1)$.
This includes a possible factor $v$, representing infinity. The height
inequality for a degree-$d$ map gives $h(\varphi(t))\ge d h(t)-O(1)$.
Applying $abc$ with exponent $1+\epsilon$, and choosing $\epsilon>0$
so that $d\epsilon/(1+\epsilon)<\delta$, consequently gives
\[
 \log\rad F(u,v)\ge(m-2-\delta)h(u/v)-C,
\]
where $F$ is the fixed integral homogeneous product of the marked
linear factors. Rational points mapping to $0,1,\infty$ but not equal
to one of the marked contacts form a fixed finite set and can be
absorbed by increasing $C$; thus no zero triple was used in $abc$.

For a nonzero integer $N$, write $\operatorname{sf}(N)$ for the product
of its odd-valuation primes. Prime by prime,
$\rad(N)^2\le |N|\operatorname{sf}(N)$.
Also $|F(u,v)|\le C'\max(|u|,v)^m$. Outside $S$ the linear factors
are pairwise coprime. A prime dividing $v$ does not divide any of them,
since $(u,v)$ is primitive. Hence their odd squarefree part outside $S$
is exactly $M(u/v)$; the part inside $S$ is bounded by the fixed product
of primes in $S$. Combining the two inequalities proves the result.
\end{proof}

Take
\[
 \Theta=\{-5,-2,-1,0,2,3,4,7\}.
\]
The fixed exceptional set used in~\cite{Parent} contains every prime
needed in Lemma~\ref{lem:marked}. With $m=8$ and $\delta=1/4$ it gives
\begin{equation}\label{eq:contact}
 h(t)\le\frac27\log M(t)+C_0\qquad(t\notin\Theta).
\end{equation}
Only ordinary rational $abc$ is asserted here. No claim that this
estimate follows from one of the weaker fixed-exponent hypotheses of
the 13-variable families is made.

\section{A parity-sensitive mixed-power estimate}
Suppose that the weak rational $abc$ inequality holds with exponent
$\kappa\in[1,4/3)$ and constant $K_\kappa\ge1$:
\[
 C\le K_\kappa\rad(ABC)^\kappa
 \quad(A+B=C,\ A,B,C\in\Z_{>0}\text{ pairwise coprime}).
\]

\begin{lemma}\label{lem:mixed}
For $x,y\in\Q$, set $s=x^8+y^2$. Then
\[
 h(x)\le \frac{2+\kappa}{4(4-3\kappa)}h(s)+C_1,
 \qquad
 h(y)\le \frac{2+\kappa}{4-3\kappa}h(s)+4C_1,
\]
where $C_1=\log K_\kappa/(2(4-3\kappa))$ suffices.
\end{lemma}
\begin{proof}
The cases $xy=0$ follow directly from the height of a power; if both
are zero all heights vanish. Otherwise write $s=M/N>0$ in lowest
terms, and take pairwise coprime positive integers $A+B=C$ such that
\[
 x^8=sA/C,\qquad y^2=sB/C.
\]
Because $A/B=(x^4/y)^2$, both $A$ and $B$ are integer squares.
We first establish the exact majorant
\begin{equation}\label{eq:radical}
 \rad(ABC)^8\le A B^4 C M^7 N^3.
\end{equation}
At a prime put $a=v_p(A)$, $b=v_p(B)$, $c=v_p(C)$ and
$e=v_p(s)$. At most one of $a,b,c$ is positive, $a,b$ are even, and
\[
 a-c+e\equiv0\pmod8,\qquad b-c+e\equiv0\pmod2.
\]
If $a>0$ and $e\ge0$, then $a+e$ is a positive multiple of eight,
so $1-a/8\le e/8$. If $a>0$ and $e<0$, then $-e\ge2$, while
$a\ge2$, and $1-a/8\le3/4\le3(-e)/8$.
If $b>0$, then $b\ge2$ and the contribution $b/2$ already covers the
radical prime. If $c>0$ and $e=0$, then $c\ge8$. If $c>0$ and $e>0$,
$1-c/8\le7/8\le7e/8$. Finally, if $c>0$ and $e<0$, the positive
integer $c-e$ is divisible by eight, giving
$1-c/8\le(-e)/8$. These cases prove
\[
 {\boldsymbol 1}_{p\mid ABC}
 \le a/8+b/2+c/8+(7/8)\max(e,0)+(3/8)\max(-e,0),
\]
and hence~\eqref{eq:radical}. Primes outside $ABC$ cause no difficulty.

Put $H=\log C$, $m=\log M$, and $n=\log N$. Since $A,B\le C$,
\[
 \log\rad(ABC)\le\frac34H+\frac78m+\frac38n.
\]
The weak $abc$ inequality yields
\[
 H\le\frac{\kappa(7m+3n)}{2(4-3\kappa)}
        +\frac{4\log K_\kappa}{4-3\kappa}
 \le\frac{5\kappa}{4-3\kappa}h(s)
        +\frac{4\log K_\kappa}{4-3\kappa}.
\]
Finally $8h(x)\le h(s)+H$ and $2h(y)\le h(s)+H$, by the displayed
primitive normalization and the product-height inequality. Division
by eight and by two proves the assertions.
\end{proof}

\begin{lemma}\label{lem:shift}
For a nonnegative rational $s$ and a positive integer $B$,
\[
 h(s)\le h(B+s),\qquad h(B+s)\ge\log B.
\]
\end{lemma}
\begin{proof}
Write $s=M/N$ in lowest terms, $M\ge0$, $N>0$. The fraction
$(M+BN)/N$ is still reduced. Its positive numerator is at least
$M$, at least $N$, and at least $BN$.
\end{proof}

Ordinary $abc$ supplies $\kappa=101/100$. For
\begin{equation}\label{eq:newscale}
 a=B+X^8+T^2
\end{equation}
Lemmas~\ref{lem:mixed} and~\ref{lem:shift} therefore give constants
independent of $B,X,T$ such that
\begin{equation}\label{eq:mixedcoefficients}
 h(X)\le\frac{301}{388}h(a)+C_1,
 \qquad h(T)\le\frac{301}{97}h(a)+4C_1,
 \qquad h(a)\ge\log B.
\end{equation}
The square property of both primitive summands in
Lemma~\ref{lem:mixed} is essential. Dropping it gives a weaker estimate
which does not close the absorption below.

\section{The explicit quadratic source}\label{sec:source}
A field coordinate means two rational coordinates in the basis
$1,\rho$; a bar denotes conjugation. Use the following coordinates:
\[
 \begin{array}{ll}
 \text{input, slope:}&X,\lambda\in\Q;\\
 \text{reference:}& f_H,S,T_H,F_H,G_H\in K;\\
 \text{point pair:}&f,h,g,k,s,F_U,F_V\in K;\\
 \text{ratios and comparison squares:}&L,L_0,C_2,I_2\in K;\\
 \text{height coordinates:}&T,z_2,z_4\in\Q;\\
 \text{second argument:}&t_1\in\Q.
 \end{array}
\]
For each of the two detector blocks retain fourteen rational scalar
coordinates $R,Z,w,w_2,w_3,\eta,b_0,e,u,p,v,q,p_2,q_2$ and three
field coordinates $A_1,A_2,\Omega$. These are twenty coordinates per
block. The total is
\[
 2+10+14+8+3+1+2\cdot20=78.
\]

\paragraph{Reference and pair equations.}
Use the same four field quadrics for $(f_H,S,F_H,G_H,T_H)$ as in the
complete reference chart of~\cite{Parent,Reference}. Their fixed
rational coefficients are unchanged; the supplied reference source is
byte-pinned. They recover the allowed finite points $H,H-P$ and the
specified reciprocal squares, including the exceptional cases proved
there. For the point pair use
\begin{equation}\label{eq:pair}
 F_U=f^2,\quad F_V=g^2,\quad h^2=f-DfF_U,\quad
 s(f-g)=h+k,\quad s(h-k)=1-D(F_U+fg+F_V).
\end{equation}
These are five field quadrics. They describe $U,V\in E(K)$, including
identity numerator points, with $U\ne V$ and finite difference $A=U-V$.
In particular
\[
 x(A)=s^2+D(f+g),\qquad y(A)=Dh-s(Df+x(A)).
\]
The source proof includes $f=g$ (tangency), torsion differences, and
all denominator exceptions. Use the same four field ratio equations
\begin{equation}\label{eq:ratios}
 hL+fS=0,\quad(1-DF_U)L+hS=0,\quad
 kL_0+gT_H=0,\quad(1-DF_V)L_0+kT_H=0.
\end{equation}

\paragraph{New input and height equations.}
Set, as linear expressions rather than additional coordinates,
\[
 \beta=z_2+3,\qquad C=L-\overline{L_0},\qquad I=L-\beta.
\]
Impose
\begin{equation}\label{eq:newrows}
 \begin{gathered}
 C_2=C^2,\qquad I_2=I^2,\qquad
 T=2\operatorname{Re}(s^2+D(f+g)),\\
 z_2=X^2,\qquad z_4=z_2^2,\qquad
 t_1-\lambda=B+z_4^2+T^2.
 \end{gathered}
\end{equation}
Thus $\beta=X^2+3$ and $t_1-\lambda=a$ from~\eqref{eq:newscale}.
Every equation is quadratic. Compared with the 79-variable source,
the extra real comparison coordinate and its defining row are deleted;
the two already retained height coordinates are reused for $X^2,X^4$.
No root, inverse, or inequality witness is silently omitted.

\paragraph{Unchanged detector and pole equations.}
For $t=t_0=\lambda$ and $t=t_1$, use independent copies of
\begin{align}
 R&=t^2-2t, & RZ+Z-R^2+35R&=0, & (R+1)w&=Z^2+24Z,\notag\\
 w_2&=w^2, & w_3&=ww_2,\notag\\
 w_3A_1&=w_2^2+w_3C_2+w_2F_H+F_U^2,\notag\\
 w_3A_2&=w_2^2+w_3I_2+w_2\overline{G_H}+\overline{F_V}^{\,2},\notag\\
 (R+1)\Omega&=A_1A_2,\qquad
 \omega_0^2-2\omega_0\eta-2\omega_1^2=0,\qquad w_3b_0=\eta^2,\notag\\
 p_2&=p^2,\qquad e^2+b_0p_2=1,\qquad
 uv=e(1-b_0),\qquad q_2=q^2,\qquad u^2-v^2=b_0q_2.
 \label{eq:detectors}
\end{align}
Here $\Omega=\omega_0+\rho\omega_1$, and rational terms in field
equations are embedded in $K$. Each block has eighteen rational rows.
The reference, pair, ratio and square rows give thirty rational rows;
the trace and two power rows give three more; the argument row gives
one, and the two detector blocks give thirty-six. Hence there are
exactly seventy rows.

Clear rational coefficients row by row to primitive integer rows
$q_{B,j}$ and define $Q_B=\sum_{j=1}^{70}q_{B,j}^2$.
Over $\Q$ its zero set is exactly the joint zero set of the rows.
The input $X$ occurs only in $z_2-X^2=0$. That row is already primitive;
its square has coefficient one at $X^4$, and no other row contains $X$.
Therefore $Q_B$ has degree exactly four and integer content one for
all fixed integer $B$. This proves those assertions without relying
on an expansion at selected numerical parameter values.

\section{Soundness with a uniform sufficient shift}
All reference, pair, ratio, contact, and pole equations in the preceding
section are unchanged from the source. In particular their domain proof
gives $t_j\notin\Theta$: writing $r=R+1$, $z=Z+37$ yields
$r=(t-1)^2$ and $r(z-r)=36$, so $r\ne0$. If $w=0$, the two factor
equations force $f=g=0$, and the pair equations yield $0=1$.
No division by a potentially zero new coordinate has been introduced.

Write the index of $A$ as $n\in\tfrac12\Z$, and set $\ell=2n\in\Z$.
The unchanged local extraction implication in~\cite{Parent} says that,
at each prime counted by $M(t_j)$, $n\equiv\beta\pmod p$ and $\beta$
is $p$-integral. It applies to the present $I=L-\beta$ verbatim: the
proof uses the value of the target but not its formula in $X$ or $B$.
The source all-point trace theorem gives fixed constants $A_0>0,J\ge0$
with
\begin{equation}\label{eq:trace}
 \ell^2\le A_0 h(T)+J.
\end{equation}
This includes torsion index zero; the additive constant is ineffective.

Suppose $n\ne\beta$. Write $X=u/v$ in lowest terms, $v>0$, and put
\[
 N=|\ell v^2-2u^2-6v^2|>0.
\]
At the consumed odd primes, $p$-integrality of $\beta=X^2+3$ implies
$p\nmid v$. The local congruences therefore give
$M(t_j)\mid N$ separately for $j=0,1$. No coprimality between these
two radicals is assumed. Moreover
\begin{equation}\label{eq:N}
 N\le (|\ell|+8)e^{2h(X)}
 \le8(1+|\ell|)e^{2h(X)}.
\end{equation}
Using~\eqref{eq:contact} and $a=t_1-t_0$ gives
\[
 h(a)\le h(t_0)+h(t_1)+\log2\le\frac47\log N+C_2.
\]
Combine this with~\eqref{eq:mixedcoefficients}. Then
\[
 h(X)\le\frac{301}{679}\log N+C_3
 \le\frac{602}{679}h(X)+\frac{301}{679}\log(1+|\ell|)+C_4.
\]
Since $602/679<1$, subtraction yields
\[
 h(X)\le\frac{301}{77}\log(1+|\ell|)+C_5.
\]
The companion inequality for $T$, together with~\eqref{eq:N}, gives
\[
 h(T)\le\frac{1204}{77}\log(1+|\ell|)+C_6.
\]
All constants are independent of $B$ and every witness. Equation
\eqref{eq:trace} now bounds $|\ell|$ uniformly, because a quadratic
function eventually exceeds any fixed multiple of $\log(1+|\ell|)$.
It follows successively that $h(X)$, $\log N$, and $h(a)$ are bounded
by fixed constants. But $h(a)\ge\log B$. For every sufficiently large
fixed odd $B$, this is impossible.

Consequently $n=X^2+3$. Hence $2X^2=\ell-6$ is an integer. In a reduced
fraction $X=u/v$ this forces $v^2\mid2$; thus $v=1$ and $X\in\Z$.
This proves soundness, including every exceptional point stratum
allowed by the unchanged source equations.

\section{Completeness for every integer, including zero}
Fix an integer $X$ and an odd $B\ge81$. Set
\[
 \beta=X^2+3\in\Z_{\ge3},\qquad A=[\beta]P,
 \qquad T=\Tr x(A),\qquad a=B+X^8+T^2>12.
\]
These are fixed before choosing any comparison prime. The source
completeness construction applies to any positive integer target
$\beta\ge3$ and any fixed positive rational additive scale $a>12$.
For clarity, its two contact lists here are $\Theta$ and $\Theta-a$.
They are disjoint, since the difference of two contacts has absolute
value at most twelve. Their rational denominators and all finite
preparation data are fixed before the retained simultaneous-prime
application. Nothing in that application requires $a$ to be a sum of
fourth powers.

Choose the source compact real neighborhoods and finite progression for
$H=[h]P$. The same prime theorem and compatible order congruences then
give the comparison slope $\lambda$ and such an $h$ with
\[
 U=[h\beta]P,\qquad V=[(h-1)\beta]P,
 \qquad U-V=A.
\]
The unchanged formal-logarithm estimates make the two exact ratios
$\beta$ to the prescribed local precision; hence the comparison tests
$C$ and $I=L-\beta$ have precisely the orders used by the source.
The fixed real neighborhoods avoid its forbidden point and denominator
values. Since the scale is fixed, letting $\lambda$ increase makes
both $\lambda$ and $\lambda+a$ tend to infinity, exactly as before.
The source finite and real preparation therefore supplies each actual
pole input in $(0,1)$ with square reduced denominator, a square in
$\Q_2$, and $2$-adic valuation at least three. The same simultaneous-twist
pole realization supplies independent pole witnesses. No equality or
square-class identification between the two inputs is imposed.

Finally set $z_2=X^2$, $z_4=X^4$. Every equation
\eqref{eq:newrows} is then satisfied, and all other rows are the source
rows just realized. This proves completeness for every integer $X$ and
all fixed odd $B\ge81$. In particular $X=0$ uses $\beta=3$, so no extra
factor of $X$, disjunction variable, or change of degree is needed.
Combining completeness with soundness proves Theorem~\ref{thm:main}.

\section{Verification scope and other routes}
The independent checker binds the exact parent sources, maps all
coordinates by name, preserves every unchanged row, derives the five
changed rows separately, and expands all seventy squares. It checks
unit content for every parameter using the input fourth power,
coefficient profiles, isolated reconstruction, and deliberate incorrect
height and input substitutions. Further exact tests cover the primewise
majorant~\eqref{eq:radical}, rational examples, and the strict fraction
$602/679$. These are finite checks, not independent proofs of the
infinite arithmetic inputs.

The mixed-power lemma can also be used as a separate height estimate.
Its application here uses the stronger eight-contact consequence of
ordinary $abc$, so it does not by itself improve the unconditional
25-variable source or either fixed weak-$abc$ 13-variable count.
The 200-variable inherited-premise quartic is left unchanged. No global
minimum in the number of variables is asserted.

\begin{thebibliography}{9}
\bibitem{Parent}\emph{A 79-variable quartic under rational $abc$},
project manuscript, corrected edition 14 November 2026,
\nolinkurl{papers/quartic_existence/quartic_abc79_20261114_v27.tex}.
\bibitem{Reference}\emph{A complete compressed reference-point chart},
project manuscript and exact coefficient source,
\nolinkurl{papers/elliptic_charts/compressed_reference_20261114_v27.tex}.
\bibitem{Pole}\emph{Pole equations from simultaneous congruent-number twists},
project manuscript, 14 November 2026,
\nolinkurl{papers/pole_parameters/congruent_twist_poles_20261114_v27.tex}.
\bibitem{Jaswal}M.~Jaswal, \emph{A proof that $\Z$ is Diophantine in $\Q$
through congruent-number twists}, Hexagon 2611.00331v2, Theorem 1.2.
The theorem is retained with exactly the preparation scope of~\cite{Pole}.
\bibitem{Elkies}N.~D.~Elkies, \emph{ABC implies Mordell},
International Mathematics Research Notices \textbf{1991}, no.~7, 99--109.
The marked Belyi argument used here is supplied explicitly in
Lemma~\ref{lem:marked}; no quantitative or effective Belyi bound is used.
\end{thebibliography}
\end{document}
